\section{Tool stability and normalization}
\begin{definition}[Tool-stable invariant regime]\label{def:tool-stable-invariant-regime}
A compact paired conservative normal form lies in a tool-stable invariant regime if monotone selectors are compression-stable, close-side states are route-faithfully robust, promote-side continuation is honest rather than a pure weight artifact, and the preserved route-faithful invariants remain unchanged under admissible uses of the tool layer.
\end{definition}

\begin{definition}[Canonical normalization regime]\label{def:canonical-normalization-regime}
A canonical normalization regime is an admissible class of normalization scale triples in which the selector class and close/promote classification remain unchanged under bounded side-preserving scale replacements.
\end{definition}

\begin{theorem}[Normalization theorem]\label{thm:normalization-theorem}
Every compact paired conservative normal form equipped with an admissible normalization scale triple admits a normalized live residual reading and a normalized weighted residual energy
\[
\Ew=w_{\mathrm{CL}}|\widetilde{\rho}_{\mathrm{CL}}|^2+w_{\Theta}|\widetilde{\rho}_{\Theta}|^2+w_A|\widetilde{\rho}_{A}|^2.
\]
Inside a canonical normalization regime, comparison, selection, and close/promote testing are independent of arbitrary raw scale imbalance.
\end{theorem}

\begin{proof}
Normalization is performed channel-wise inside the route-faithful split. Within the canonical normalization regime, equivalent admissible scale choices preserve the relevant comparisons and branch classifications, so the normalized state is the preferred pre-OP working surface.
\end{proof}

